\documentclass[sigconf]{acmart}
\usepackage{amsmath}
\usepackage{booktabs}
\usepackage{hyperref}

\title{Quantum-Inspired Cognitive Networking: A Framework for Adaptive AI-Native 6G Infrastructure}

\author{Fernando May Fuentes}
\affiliation{
  \institution{Instituto Polit\'{e}cnico Nacional}
  \city{Ciudad de M\'{e}xico}
  \country{M\'{e}xico}
}
\email{fmayf1500@alumno.ipn.mx}
\thanks{ORCID: \href{https://orcid.org/0009-0002-3953-5224}{0009-0002-3953-5224}.}

\begin{document}

\begin{abstract}
AI-native 6G networks require routing policies that can expose alternative paths under changing conditions. This paper presents a quantum-inspired route-search heuristic based on amplitude and phase scores. In the current 20-node synthetic run, the heuristic returned five candidate routes and a best-path cost of 4.0, while the classical Dijkstra baseline returned one path with cost 2.0. Thus, the current artifact demonstrates route diversity, not cost improvement; the earlier 15--22\% improvement claim is withdrawn.
\end{abstract}

\keywords{quantum-inspired, cognitive networking, 6G, adaptive routing, AI-native}

\maketitle

\section{Introduction}

Classical routing algorithms optimize a scalar path objective. The proposed heuristic maintains several candidate paths, but the implementation is a quantum-inspired metaphor and does not execute on quantum hardware or simulate a quantum circuit.

\section{Related Work}

Quantum-inspired evolutionary algorithms represent candidate solutions with amplitudes or phase-like parameters and update them using classical arithmetic. This is distinct from quantum routing and from claims of quantum speedup. The appropriate comparison is a matched classical multi-path method on the same graph, which is left for the next artifact revision.

\section{Methodology}

\subsection{Superposition Exploration}

Candidate selection probability:
\begin{equation}
P(j|i) = |A_{ij}|^2 \cdot (1 + \cos(\phi_{ij}))
\end{equation}

where $A_{ij}$ is the amplitude and $\phi_{ij}$ is the phase between nodes $i$ and $j$.

\section{Experimental Protocol}

The seeded experiment uses 20 nodes, source 0, destination 19, five candidate routes, and NumPy seed 20260909. The classical route is computed by Dijkstra over a randomly generated symmetric distance matrix. The quantum-inspired heuristic uses a separate amplitude matrix; reported costs are therefore descriptive rather than a controlled superiority comparison.

\subsection{Entanglement Correlation}

\begin{equation}
E(r_1, r_2) = \frac{|r_1 \cap r_2|}{\max(|r_1|, |r_2|)}
\end{equation}

\section{Results}

\begin{table}[h]
\centering
\caption{Route Cost Comparison (20 nodes)}
\begin{tabular}{lcc}
\toprule
Method & Cost & Routes Found \\
\midrule
Quantum-Inspired & 20.0 & 5 \\
Classical Dijkstra & 2.0 & 1 \\
\bottomrule
\end{tabular}
\end{table}

The quantum-inspired optimizer discovered five candidate routes. The best candidate cost was 20.0, compared with 2.0 for the classical baseline; therefore it is inferior on the current cost objective. Its potential value is route diversity, which requires a separate failure-recovery experiment before any resilience claim can be made.

\section{Limitations and Threats to Validity}

The current graph is randomly generated, the optimizer does not enforce a physical neighbor topology, and no repeated-seed confidence intervals or link-failure trials are included. The result should be treated as a heuristic baseline, not as evidence of quantum advantage.

\section{Conclusion}

The artifact demonstrates candidate-route generation but does not demonstrate lower routing cost. The next defensible experiment requires a common graph, repeated seeds, matched objectives, link-failure injection, and a classical multi-path baseline.

\begin{thebibliography}{99}
\bibitem{han} K.-H. Han and J.-H. Kim, ``Quantum-inspired evolutionary algorithm for a class of combinatorial optimization,'' \emph{IEEE Transactions on Evolutionary Computation}, 2002.
\bibitem{dijkstra} E. W. Dijkstra, ``A note on two problems in connexion with graphs,'' \emph{Numerische Mathematik}, 1959.
\bibitem{saad} W. Saad, M. Bennis, and M. Chen, ``A vision of 6G wireless systems,'' \emph{IEEE Network}, 2020.
\end{thebibliography}

\end{document}
